\documentclass{article}
\title{ECE 496 -- Homework 1}
\usepackage{graphicx}
\usepackage{pdfpages}
\usepackage{fancyhdr}
\usepackage{enumitem}
\usepackage{pdfpages}
\usepackage{amsmath}
\usepackage{amssymb}
\pagestyle{fancy}
\fancyhf{}
\fancyhead[L]{Gianni Avilan-Losee}
\fancyhead[C]{ECE 496}
\fancyhead[R]{Homework 1}
\fancyfoot[C]{\thepage}
\usepackage{microtype}
\usepackage{textcomp, gensymb}
\author{Gianni Avilan-Losee}
\setlist[enumerate, 1]{label=\thesection.\arabic*}
\setcounter{section}{1}
\begin{document}
\begin{enumerate}
	\setcounter{enumi}{1}
	\item
	\begin{enumerate}[label=\alph*)]
		\item 
			\begin{tabular}{|c|c|}
				\hline
				Diameter (mm) & Area (cm\(^2\)) \\
				\hline
				25 & 4.91 \\
				\hline 
				50 & 19.64\\
				\hline
				75 & 44.18\\
				\hline
				100 & 78.54 \\
				\hline
				125  & 122.72\\
				\hline
				150 & 176.71\\
				\hline
				200 & 314.16\\
				\hline
				300 & 706.86\\
				\hline
				450 & 1590.43\\
				\hline
			\end{tabular} 
     \\
	 \item Dividing wafer area by die area, approximately 159,043 \(1\text{mm} \times 1\text{mm}\) dice.
		 \\
	 \item Dividing wafer area by die area, \[
			 \frac{159,043.13\text{mm}^2}{625\text{mm}^2} = 254.47,
	 \] 
	 so approximately 254 dice.
	\end{enumerate}
	\setcounter{enumi}{4}
\item Substituting \(Y = 2020\) into the given formula, \[
		N(2020) = 1027 \times 10^{0.1505(2020-1970)} = 3.44 \times 10^{10},
\]
around 34 billion transistors per microprocessor.
\setcounter{enumi}{6}
		\begin{enumerate}[label=\alph*)]
\item Using the given formula, \[
		N(Y) = 2N(Y) \Rightarrow Y = \frac{\log_{10}(2)}{0.1505} = 2.002, 
\] 
and thus the transistor count on a microprocessor would be expected to double about every two years. 
\item Using the same formula, \[
		Y = \frac{\log_{10}(10)}{0.1505} = 6.6445,
\] and thus the transistor count on a microprocessor would be expected to increase by a factor of 10 roughly every six and half years.
\end{enumerate}
\setcounter{enumi}{11}
\item
\begin{enumerate}[label=\alph*)]
	\item Using Fig 1.1(b), we may assume that approximately 175 \(10 \times 10\text{mm}\) dice are on a \(150\text{mm}\) wafer. Assuming 35\% of the dice are good, this leaves us with \(175*0.35 = 61.25\), so approximately 61 good dice. Assuming a cost per wafer of \$1,000, the cost per die is approximately \[
	 \frac{\$1000}{61} = \$16.39.
	\] 
\item Again using Fig 1.1(b), we may assume that approximately 314 \(10 \times 10\text{mm}\) dice are on a \(200\text{mm}\) wafer. Using the same perecentage of good dice as above, this leaves \(314*0.35 = 109.90\), or approximately 109 good dice. Then, approximate cost per wafer is \[
  \frac{\$1000}{109} = \$9.17.
\] 
\end{enumerate}
\setcounter{enumi}{13}
\item See below for included flowchart:
	\includepdf{homework-1_cover-page.pdf}
\end{enumerate}
\end{document}
